%% =====================================================================
%%  SUPPLEMENTARY MATERIAL
%%  "Efficient and practical solutions for the correlated storage
%%   location assignment problem in automated pick-and-pass warehouses"
%%  International Journal of Production Research
%%
%%  This file holds the calibration and sensitivity study of the
%%  clustering heuristic, referenced from Section 4.1 of the main
%%  article. It compiles standalone: it needs only IJPR_CSLAP.bib
%%  alongside it. Cross-references to the main article are written out
%%  in full, since the two documents compile separately.
%% =====================================================================
\PassOptionsToPackage{table}{xcolor}
\documentclass[11pt]{article}

\usepackage[a4paper,margin=2.5cm]{geometry}
\usepackage{amsmath,amssymb,amsfonts}
\usepackage{mathtools}
\usepackage{booktabs}
\usepackage{float}
\usepackage{xcolor}
\usepackage{tabularx}
\usepackage{multirow}
\usepackage{algorithmicx}
\usepackage{algorithm}
\usepackage{algpseudocode}
\usepackage{url}
\usepackage{hyperref}

\usepackage[natbibapa,nodoi]{apacite}
\setlength\bibhang{12pt}

% Number sections, tables and figures with an S prefix, as T&F expect
% for supplementary items.
\renewcommand{\thesection}{S\arabic{section}}
\renewcommand{\thetable}{S\arabic{table}}
\renewcommand{\thefigure}{S\arabic{figure}}
\renewcommand{\thealgorithm}{S\arabic{algorithm}}

% The main article uses \tbl{} and \tabnote{} from interact.cls; supply
% equivalents so the table bodies below are unchanged from the article.
\newcommand{\tbl}[1]{\caption{#1}}
\newcommand{\tabnote}[1]{\vspace{2pt}\footnotesize\raggedright #1\par}

\title{Supplementary material:\\
Calibration, proofs and implementation record}
\author{Ermal Belul \and Marwane Bouznif \and Dritan Nace \and Antoine Jouglet}
\date{}

\begin{document}
\maketitle

\noindent
This supplement holds the material the main article references but does not carry.
Sections~\ref{sup:capsweep} to~\ref{sup:heuristic} evaluate the design parameters of the greedy
clustering heuristic of Section~4.1 of the main article: whether the community size bound $\beta$
scales with catalogue size ($N$) or with station slot capacity ($\zeta$), how sensitive the
solution is to each threshold ($\phi, \gamma, r, \beta$), the paired statistical tests behind the
synthetic benchmark, and the pseudocode of the procedure. Section~\ref{sup:impl} gives the
implementation, hardware and parameter record for every method; Section~\ref{sup:proofs} proves
the two bounding claims of the main article; Section~\ref{sup:betasearch} reports the full
industrial $\beta$ search; Section~\ref{sup:cgrecord} describes the released per-instance
column-generation convergence record; and Section~\ref{sup:relocap} gives the column-generation
formulation of the bounded-reassignment cap.

To isolate these effects without confounding catalogue volume and station capacity (which are
tied in the main 29-instance benchmark of Section~5 of the main article), we generated an
expanded dataset of 80 synthetic instances across crossed combinations of $N \in \{500, 1000\}$
and $\zeta \in \{20, 25, 50, 100\}$, using the same common-itemset generator script as the
primary benchmark but with independent random seeds. This secondary dataset allows us to
systematically test whether $\beta$ scales with $N$ or $\zeta$, and to test threshold
robustness across diverse warehouse configurations.

\section{Community bound scaling across warehouse geometries}\label{sup:capsweep}
The three filters of Section~4.1 of the main article are indexed on the catalogue and the
community bound on the station. That asymmetry needs evidence, because the benchmark of
Section~5 cannot supply it: its generator ties the station count to the catalogue, so
$\zeta=100$ at every size from 500 SKUs upward and the ratio $\beta/\zeta$ is constant across
the whole family. Catalogue size and slot capacity are confounded there, and no result measured
on it can separate them.

We therefore built a second family in which the two are not confounded. It uses the same
generator as Section~5 of the main article (the common-itemset procedure of
\citet{zhang2019cie} at correlation degree $0.7$, with order size drawn from $U[5,15]$ and
proportional order volumes) and changes one thing: the number of stations is set independently
of the catalogue rather than as $\max(5,\lfloor N/100\rfloor)$. Two catalogue sizes,
$N\in\{500,1000\}$, are crossed with four station counts chosen so that the same slot
capacities $\zeta\in\{20,25,50,100\}$ occur at both sizes. Ten instances per cell give 80 in
all, generated on seeds $2001$ to $2010$, disjoint from the seeds of the benchmark. Station
labels are drawn from a separate random stream, so at a fixed seed the order structure is
identical across the four geometries and only the station configuration changes. Any difference
between cells is therefore attributable to geometry alone.

Each instance is solved at every $\beta$ on the grid
$\{2,3,5,8,10,12,15,18,20,25,30,40,50,75,100\}$ that falls at or below its own $\zeta$, which
gives 980 runs. All 980 returned a feasible layout. Table~\ref{tab:capsweep} reports the
minimising bound in each cell. Both instance families, with the generator settings and seeds
needed to reconstruct them exactly, accompany the article.

\begin{table}[H]
\centering
\tbl{Visit-minimising community bound $\beta^\ast$ across warehouse geometries, on 80 instances disjoint from the benchmark. The same $\zeta$ appears at both catalogue sizes.}
{\small\begin{tabular*}{\textwidth}{@{\extracolsep{\fill}}rrrrrr@{}}
\toprule
$\zeta$ & $N$ & $|S|$ & $\beta^\ast$ & $\beta^\ast/\zeta$ & $\Delta$ visits vs.\ $\beta=15$ \\
\midrule
20  & 500   & 25 & 12 & 0.60 & $-0.72$\% \\
20  & 1{,}000 & 50 & 12 & 0.60 & $-0.45$\% \\
25  & 500   & 20 & 15 & 0.60 & $0.00$\% \\
25  & 1{,}000 & 40 & 15 & 0.60 & $0.00$\% \\
50  & 500   & 10 & 30 & 0.60 & $-0.87$\% \\
50  & 1{,}000 & 20 & 30 & 0.60 & $-0.84$\%\textsuperscript{$\ast$} \\
100 & 500   & 5  & 40 & 0.40 & $-1.01$\%\textsuperscript{$\ast$} \\
100 & 1{,}000 & 10 & 40 & 0.40 & $-1.06$\%\textsuperscript{$\ast$} \\
\bottomrule
\end{tabular*}}
\tabnote{$\beta^\ast$ minimises mean visits in its cell over the swept values, which include 15. The last column is the change from 15 to $\beta^\ast$, so a negative entry means $\beta^\ast$ uses fewer visits and a zero entry means 15 is itself the minimiser, as it is in both $\zeta=25$ cells. \textsuperscript{$\ast$}Separates from 15 at $p<0.05$ under a Holm-corrected Wilcoxon signed-rank test over the ten instances of the cell. The unstarred non-zero cells fail that test on within-cell variance rather than on effect size: at $\zeta=50$, $N=500$ the point effect of $-0.87$\% exceeds that of the significant cell beneath it.}
\label{tab:capsweep}
\end{table}

At a fixed slot capacity $\beta^\ast$ is the same at both catalogue sizes in all four cells,
whereas at a fixed catalogue size it moves with $\zeta$ across the whole range. Over the eight
cells the rank correlation between $\beta^\ast$ and $\zeta$ is $0.962$, and between
$\beta^\ast$ and $N$ it is $0.000$. This confirms that the bound scales with the slot capacity
of a station and is insensitive to the length of the catalogue.

The dependence is monotone but not proportional. The minimising ratio $\beta^\ast/\zeta$ is
$0.60$ at $\zeta\le 50$ and falls to $0.40$ at $\zeta=100$, so a single coefficient cannot be
optimal everywhere. We adopt $0.4$ for two reasons. It is the minimising value at $\zeta=100$,
the capacity of every benchmark instance of 500 SKUs and above and of the sizes at which the
method is intended to be used. Where it is not the minimiser it errs low, which matters because
larger communities are harder to place within the workload constraint, as Section~6 of the main
article shows on a site whose stations are not interchangeable. The price of that choice is nil
at $\zeta=100$, where $0.4\zeta$ is the minimiser, and between $0.15$ and $1.03$\% of visits in
the six cells below it.

\section{Threshold sensitivity of the clustering heuristic}\label{sup:sensitivity}
Table~\ref{tab:sensitivity} rescales each of the four thresholds of Section~4.1 of the main
article from half to twice its nominal value, one at a time and then all four together, on the
twelve 50-SKU and ten 500-SKU instances. The 1,000- and 2,000-SKU instances are omitted because
four and three instances are too few for an informative interval, although the ordering of the
effects there agrees with what follows. Multipliers apply to the effective nominal value on the
instance, after the floor of 5 (for the three filters that value is set by $N$ and for the
community bound by $\zeta$). At 50 SKUs both frequency floors already sit on their clamp floor
of 2, and the community bound sits on its floor of 5, so only the upward settings are distinct
there. Entries are paired mean changes against the nominal setting on the same instance. The
procedure is deterministic: repeating the nominal setting under five hash seeds returns an
identical layout on every instance, so any difference is a real threshold effect and not
run-to-run variation. Reporting follows the conventions for computational experiments with
heuristics \citep{barr1995designing, hooker1995testing}.

The two frequency floors are inert, as rescaling either leaves the visit count unchanged and
the surviving edge set bit-identical at every size, because neither removes a product or pair
that the kept-edge ratio does not already remove. The ratio changes the edge set substantially
but the objective very little, separating from the nominal setting only at $\times2$. The
community bound is the one threshold whose response is not monotone. Halving it costs $0.63\%$
at 500 SKUs and doubling it costs $1.19\%$, while $\times0.75$ and $\times1.5$ are
indistinguishable from the nominal setting, while at 1,000 and 2,000 SKUs the pattern is the
same, with doubling costing $1.84\%$ and $2.27\%$. The exception is $\times1.5$ at 500 SKUs,
which recovers $0.45\%$, consistent with Section~\ref{sup:capsweep} above, where the minimising
ratio at $\zeta=100$ is $0.40$ but rises to $0.60$ at smaller capacities, so a bound half again
as large is not yet past the minimum on every instance. A default whose response rises on both
sides at the two larger sizes, and is flat rather than falling at the smaller one, is indexed on
the right quantity, whereas a default with a monotone response is not. Leaving all four at their
defaults is therefore supported here, not assumed.

\begin{table}[H]
\centering
\tbl{Sensitivity of the clustering heuristic to its four thresholds, on the 22 instances of 50 and 500 SKUs.}
{\small\setlength{\tabcolsep}{4pt}\begin{tabularx}{\textwidth}{@{}X c c@{}}
\toprule
Threshold setting & $\Delta$ visits, 50 SKUs & $\Delta$ visits, 500 SKUs \\
\midrule
Frequency floor $\phi$, $\times0.5$ to $\times2$\textsuperscript{a}     & $0.00$\% & $0.00$\% \\
Co-occurrence floor $\gamma$, $\times0.5$ to $\times2$\textsuperscript{a} & $0.00$\% & $0.00$\% \\
\midrule
Kept-edge ratio $r$, $\times0.5$   & $-0.77$\% [$-2.13$, $+0.58$] & $+0.16$\% [$-0.09$, $+0.41$] \\
Kept-edge ratio $r$, $\times0.75$  & $-0.66$\% [$-1.98$, $+0.65$] & $0.00$\% \\
Kept-edge ratio $r$, $\times1.5$   & $+0.23$\% [$-0.72$, $+1.17$] & $+0.08$\% [$-0.27$, $+0.43$] \\
Kept-edge ratio $r$, $\times2$     & $-0.74$\% [$-2.26$, $+0.78$] & $+1.08$\% [$+0.18$, $+1.98$] \\
\midrule
Community bound $\beta$, $\times0.5$  & $+1.77$\% [$-0.17$, $+3.72$] & $+0.63$\% [$-0.07$, $+1.32$] \\
Community bound $\beta$, $\times0.75$ & $+0.25$\% [$-1.19$, $+1.69$] & $+0.03$\% [$-0.46$, $+0.51$] \\
Community bound $\beta$, $\times1.5$  & $-0.25$\% [$-2.29$, $+1.79$] & $-0.45$\% [$-1.20$, $+0.30$] \\
Community bound $\beta$, $\times2$    & $-0.25$\% [$-1.60$, $+1.10$] & $+1.19$\% [$+0.25$, $+2.13$] \\
\midrule
All four, $\times0.5$ & $+1.15$\% [$-1.14$, $+3.43$] & $+0.59$\% [$-0.08$, $+1.27$] \\
All four, $\times2$   & $-1.10$\% [$-2.57$, $+0.38$] & $+1.78$\% [$+1.02$, $+2.54$] \\
\bottomrule
\end{tabularx}}
\tabnote{Entries are mean paired changes in station visits against the nominal setting, with 95\% confidence intervals; a positive value means more visits. Every setting returned a feasible layout on every instance. \textsuperscript{a}All four settings give bit-identical output, so the rows are collapsed into one.}
\label{tab:sensitivity}
\end{table}

\section{Paired statistical testing of the synthetic benchmark}\label{sup:tests}
Section~5 of the main article compares the methods on mean station visits, gaps, running time
and workload feasibility across the 29 synthetic instances. This section reports the paired
statistical testing behind those comparisons, which is kept out of the main text because its
resolution is bounded by the instance counts and adds little to the readings the main tables
already support.

Instances differ by two orders of magnitude in difficulty, so comparing two methods on their
averages would confound the between-method difference with the between-instance variance. We
therefore compare them instance by instance. For each pair of methods we take the per-instance
difference in station visits and test whether those differences are systematically signed. The
test is the two-sided exact Wilcoxon signed-rank test \citep{wilcoxon1945}, the standard choice
for comparing algorithms across a set of benchmark instances because it imposes no
distributional assumption on the differences \citep{demsar2006}, and ties are discarded. The
reported $p$ is the probability, under the null hypothesis that the paired differences are
symmetric about zero, of observing a signed-rank statistic at least as extreme as the one
obtained, in either direction. Table~\ref{tab:tests} reports the outcome against the
set-variable MILP, the per-instance reference of the main article.

\begin{table}[H]
\centering
\tbl{Paired Wilcoxon signed-rank tests on the 29 synthetic instances, against the set-variable MILP.}
{\small\begin{tabular*}{\textwidth}{@{\extracolsep{\fill}}llrrrr@{}}
\toprule
Comparison & Size (SKUs) & $n$ & Wins & $p$ & Min.\ $p$ \\
\midrule
Binary MILP vs.\ reference & 50      & 12 & 4  & 0.74   & 0.0005 \\
                           & 500     & 10 & 8  & 0.049  & 0.0020 \\
                           & 1{,}000 & 4  & 1  & 0.38   & 0.125  \\
                           & 2{,}000 & 3  & 0  & 0.25   & 0.25   \\
                           & pooled  & 29 & 13 & 0.91   & --     \\
\midrule
CG vs.\ reference          & 50      & 12 & 1  & 0.0010 & 0.0005 \\
                           & 500     & 10 & 2  & 0.0273 & 0.0020 \\
                           & 1{,}000 & 4  & 2  & 1.00   & 0.125  \\
                           & 2{,}000 & 3  & 3  & 0.25   & 0.25   \\
                           & pooled  & 29 & 8  & 0.18   & --     \\
\midrule
GA and SA-C vs.\ reference & 50      & 12 & 0  & 0.0005 & 0.0005 \\
                           & 500     & 10 & 0  & 0.0020 & 0.0020 \\
                           & 1{,}000 & 4  & 0  & 0.125  & 0.125  \\
                           & 2{,}000 & 3  & 0  & 0.25   & 0.25   \\
\midrule
Heuristic vs.\ reference   & 50      & 12 & 0  & 0.0005 & 0.0005 \\
                           & 500     & 10 & 0  & 0.0020 & 0.0020 \\
                           & 1{,}000 & 4  & 0  & 0.125  & 0.125  \\
                           & 2{,}000 & 3  & \textbf{3}  & 0.25   & 0.25   \\
\bottomrule
\end{tabular*}}
\tabnote{``Wins'' counts instances on which the first-named method returns fewer visits. ``Min.\ $p$'' is the smallest two-sided $p$ attainable at that sample size, $2^{1-n}$, and where $p$ equals it the test is saturated rather than failed, so the win count carries the evidence. The GA, SA-C and Heuristic rows coincide at 50, 500 and 1{,}000 SKUs because none of the three wins a single paired instance there. The Heuristic separates from them only at 2{,}000 SKUs, where it wins all three.}
\label{tab:tests}
\end{table}

The tests saturate at the two larger sizes. The resolution of the test is bounded by the number
of instances available, and on this benchmark that bound is binding: at 1{,}000 and 2{,}000
SKUs no result, however unanimous, can reach the conventional 0.05 threshold. At those two
sizes we therefore treat the win count as the primary evidence, which is the form in which the
main article states them.

The separations are reported without correction for multiplicity. Of the seven separations
below 0.05 that the table shows, the five at $p \le 0.0020$ survive a Bonferroni, Holm or
Benjamini--Hochberg correction over the eighteen tests reported, whereas the two at
$p = 0.0273$ and $p = 0.049$ do not. A small $p$ indicates that the methods differ, whereas a
large $p$ indicates only that these instances are insufficient to separate them, which is a
weaker statement than equivalence. Effect size and consistency are therefore reported
separately: the Gap column of the main article's benchmark table quantifies the magnitude of a
difference, and the test quantifies how reliably it recurs.

\section{Pseudocode of the greedy clustering heuristic}\label{sup:heuristic}
Algorithm~\ref{alg:heuristic} states in full the five stages that Section~4.1 of the main
article describes and its Figure~2 sketches. Here $|s|$ is the product count of station $s$,
$\mathrm{cnt}_{ij}$ the co-occurrence count of a product pair, $\mathrm{cnt}_i$ the order
frequency of product $i$, $\mathrm{load}(s)=\sum_{p\in s}L_p/V_s$ the workload of station $s$,
and $\zeta_s$, $T_s$, $L_p$ and $V_s$ are as defined in Section~3 of the main article. The four
thresholds are restated in the input line so that the procedure can be read on its own. The
ratio test in Step~1 normalises by the \emph{larger} of the two endpoint frequencies, so it is
symmetric in $i$ and $j$ and reads as the confidence of the weaker association direction.

\begin{algorithm}[H]
\caption{Greedy product clustering and station assignment}
\label{alg:heuristic}
\begin{algorithmic}[1]
\State \textbf{Input:} products $P$, orders $O$, stations $S$; filters from $N=|P|$: floor $\phi=\max(2,\lfloor N/100\rfloor)$, minimum co-occurrence $\gamma=\max(2,\lfloor N/50\rfloor)$, kept-edge ratio $r=0.1$; size bound from the slot capacity $\zeta=|P|/|S|$: $\beta=\max(5,\lfloor 0.4\,\zeta+\tfrac12\rfloor)$.
\Statex
\State \textbf{Step 1: preprocessing and pair generation}
\State $P^{\ast}\gets\{\,p\in P:\mathrm{freq}(p)\ge\phi\,\}$
\ForAll{multi-item orders $o\in O$}
  \State form every unordered pair $\{p_i,p_j\}\subseteq P_o\cap P^{\ast}$
\EndFor
\State for each pair, $\mathrm{cnt}_{ij}\gets\lvert\{\,o\in O:\{p_i,p_j\}\subseteq P_o\,\}\rvert$; keep pairs with $\mathrm{cnt}_{ij}\ge\gamma$
\State for each kept pair, $\eta_{ij}\gets\mathrm{cnt}_{ij}/\max(\mathrm{cnt}_i,\mathrm{cnt}_j)$; keep pairs with $\eta_{ij}\ge r$
\State sort the kept pairs by decreasing $\mathrm{cnt}_{ij}$
\Statex
\State \textbf{Step 2: community formation}
\State $\mathcal{C}\gets\varnothing$; mark every product unassigned
\While{candidate pairs remain}
  \State take the top pair $\{p_a,p_b\}$ and open a community $c\gets\{p_a,p_b\}$
  \While{$|c|<\beta$ and an unassigned product is adjacent to $c$}
    \State $p_x\gets$ unassigned product of highest aggregate co-occurrence with $c$
    \State $c\gets c\cup\{p_x\}$; delete every remaining pair that contains $p_x$
  \EndWhile
  \State $\mathcal{C}\gets\mathcal{C}\cup\{c\}$
\EndWhile
\Statex
\State \textbf{Step 3: post-processing}
\ForAll{products $p$ with $\mathrm{freq}(p)<\phi$}
  \If{some community $c$ has room ($|c|<\beta$)}
    \State add $p$ to the community of highest co-occurrence with $p$
  \EndIf
\EndFor
\State group any still-unassigned products into new communities of size at most $\beta$
\Statex
\State \textbf{Step 4: station assignment}
\State call $s$ \emph{admissible} for a block $c$ when $|s|+|c|\le\zeta_s$ and $\mathrm{load}(s)+\sum_{p\in c}L_p/V_s\le T_s$
\State order the communities by cumulative demand, most impactful first
\ForAll{communities $c$ in that order}
  \If{some station is admissible for $c$}
    \State assign $c$ to the admissible station of highest historical preference
    \State \Comment{ties broken by the lowest resulting load $\mathrm{load}(s)/T_s$}
  \Else
    \State split $c$, heaviest product first, placing each product by the same rule
  \EndIf
\EndFor
\Statex
\State \textbf{Step 5: workload repair} \Comment{$\mathrm{aff}(x,s)=\sum_{y\in s}\mathrm{cnt}_{xy}$}
\While{some station exceeds its workload cap}
  \State $s\gets$ the station of greatest relative overload $\mathrm{load}(s)/T_s$
  \State $E\gets$ exchanges of a heaviest product $p$ of $s$ against a lightest product $q$
  \State \hskip 1.2em of another station $t$ that lighten $s$ and keep $t$ inside its cap
  \State \textbf{if} $E=\varnothing$ \textbf{then break}
  \State apply the $(p,q)\in E$ of least broken weight $\mathrm{aff}(p,s)+\mathrm{aff}(q,t)-\mathrm{aff}(p,t)-\mathrm{aff}(q,s)$
\EndWhile
\State \textbf{Output:} a product-to-station assignment respecting $\zeta_s$ and $T_s$
\end{algorithmic}
\end{algorithm}

\section{Implementation, hardware and baseline parameters}\label{sup:impl}
This section states in full the computational setup that Section~5.2 of the main article
summarises, together with the incremental swap evaluation of Section~4.2.3, the stage budgets of
the price-and-complete drive, and the parameters of the two literature baselines.

\subsection{Incremental evaluation of the swap descent}\label{sup:swapeval}
The descent of Section~4.2.3 of the main article keeps a support-count matrix
$\mathrm{cnt}[i][u]$, the number of products of station $i$ contained in support $u$. Station $i$
visits $u$ exactly when $\mathrm{cnt}[i][u]>0$. Removing a product frees the supports whose count
falls to zero and inserting it costs those whose count rises from zero. Only supports containing
the two moved products can change, so a swap is scored in time proportional to their support
degrees rather than by re-scanning the orders, and the counts are updated the same way once a swap
is accepted. Beyond about 300 products the candidate partner set is sampled to keep the cost of one
descent step in hand.

\subsection{Stage budgets of the price-and-complete drive}\label{sup:budgets}
The time cap $B$ of Section~4.2.3 of the main article is divided across the four stages. On the
homogeneous benchmarks the warm start ends at $0.08B$, the generation loop at $0.70B$, the integer
master at $0.85B$, and the final polish takes the remainder. The heterogeneous industrial run uses
a schedule matched to its longer runs: warm start $0.05B$, generation loop $0.62B$, a bound pass at
$0.74B$, integer master $0.82B$, and a closing polish.

\subsection{Baseline operators and parameters}\label{sup:baselines}
Both baselines score a candidate layout with the same surrogate: the station-visit count, plus
1{,}000 times the total excess over slot capacity or workload cap, plus 0.01 times the
co-occurrence weight of product pairs placed at different stations. The first term is the real
objective, the second makes an infeasible layout uncompetitive, and the third breaks ties between
layouts the visit count cannot separate.

The genetic algorithm carries 50 assignment vectors for at most 200 generations or until the
per-size cap of Section~5.2 of the main article is reached, with binary-tournament parent
selection, crossover at probability 0.8 copying a contiguous segment of one parent into the other
up to a quarter of the catalogue with a repair pass for over-capacity stations, mutation at
probability 0.1 exchanging two products' stations, and the best layout of each generation carried
forward unchanged. Each of the 50 initial members is drawn by assigning products to stations
uniformly at random subject to $\zeta_s$, as \citet{pan2015} prescribe.

Simulated annealing with correlated interchange starts at temperature 800 and cools geometrically
by 0.95 per level until the temperature reaches 1 or the cap is met, proposing $\min(2|P|, 1000)$
moves per level. A move is drawn from the correlated pair list, co-occurring pairs being sorted by
descending co-occurrence and split into ten equal groups from which a group and then a pair are
drawn uniformly, and the two products exchange stations. A move lowering the surrogate is always
accepted, whereas one raising it by $\Delta$ is accepted with probability $\exp(-\Delta/T)$. The
run starts from a cube-per-order-index (COI) layout, which is correlation-blind like LPT but ranks
products by demand volume per unit of storage occupied rather than by order-line frequency, filling
stations in that order \citep{kim2012}.

\subsection{Hardware, solver versions and seeds}\label{sup:hardware}
Every run reported in the main article was produced on the same remote machine, a 16-core Windows
Server 2019 host with 128\,GB of memory, under the same Python 3.10 environment, and every method
received the same core and thread allocation, four threads, so no method was given more compute
than another under the shared wall-clock cap. The column generation runs on CPLEX 22.1.1 and the
two MILP representations on Hexaly 13.0. One run per instance is reported in every cell. The two
baselines are our own Python implementations of the published operators, so a residual
implementation-efficiency difference against the two commercial engines cannot be excluded, and we
record how far each baseline got rather than assume it exhausted its search.

Both metaheuristic baselines draw their search decisions from a NumPy generator seeded at 42, so
SA-C is fully determined by the instance and the cap. The genetic algorithm's capacity repair draws
instead on the global random stream, seeded at 20240612 by the experiment harness, and that stream
is the only source of run-to-run variation in the reported GA figures. The column-generation drive
is also stochastic, since Stage 1 mutates the warm start and the descent samples candidate partners
beyond about 300 products, and it draws on the same harness stream. We report one run per cell for
it as well, so the 2,000-SKU result rests on three single draws and a replicated cell would
strengthen it.

\section{Proofs of the bounding claims}\label{sup:proofs}

\subsection{Binary MILP linear relaxation floor}\label{sup:bounds_milp}
Take any point satisfying $\sum_{s\in S} x_{ps}=1$. For an order $o \in O$ and a product
$p \in P_o$, the linking row $x_{ps}\le z_{os}$ forces $z_{os}\ge x_{ps}$ for every station, and
summing over the stations gives $\sum_{s\in S} z_{os}\ge\sum_{s\in S} x_{ps}=1$. The objective
therefore cannot fall below $|O|$ at any feasible point, fractional or integral. That floor is
attained: setting $x_{ps}=1/|S|$ meets the capacity rows whenever $|P|/|S|\le\zeta_s$ and the
workload rows under the operational slack of Section~5.1 of the main article, so $z_{os}=1/|S|$ is
feasible and the objective equals $|O|$ exactly. On this instance family the relaxation is
therefore worth precisely the statement that every order visits at least one station, independent
of the co-occurrence structure that makes one instance harder than another, so a branch-and-bound
search starts from that value with nothing instance-specific to prune with.

\subsection{Column generation Farley-type dual bound}\label{sup:bounds_cg}
The bound does not rest on the price-and-complete drive. It follows from linear-programming duality
applied to the restricted master together with a valid dual bound from the pricing solve, and would
hold whatever heuristic produced the layouts.

Section~4.2.1 of the main article aggregates the station-indexed master over interchangeable
stations. Written over the current pool $Q'$ of feasible bundles, with $c_q$ the weighted-support
cost of bundle $q$, that restricted master is
\begin{align}
\min \quad & \sum_{q \in Q'} c_q\, y_q \notag\\
\text{s.t.} \quad
& \sum_{q \ni p} y_q \ge 1 && \forall p \in P \quad (\sigma_p \ge 0) \label{sup:eq:cover}\\
& \sum_{q \in Q'} y_q \le |S| && (\mu \le 0) \label{sup:eq:card}\\
& 0 \le y_q \le 1 && \forall q \in Q'. \notag
\end{align}
The upper bounds $y_q\le 1$ may be dropped, since reducing any $y_q$ above one to one preserves
coverage, relaxes the cardinality row and cannot increase the objective, so an optimal dual
solution with zero multipliers on them exists. Strong duality gives
$z_{\mathrm{RMP}}=\sum_p\sigma_p+|S|\mu$. The pair $(\sigma,\mu)$ need not be dual-feasible for the
full master, because a bundle outside the pool may price out negatively. Let $\mathrm{rc}_{\min}$
be the least reduced cost over all feasible bundles and suppose it is negative. Replacing $\mu$ by
$\mu'=\mu+\mathrm{rc}_{\min}$ restores feasibility, since every bundle then satisfies
$c_q-\sigma(q)-\mu'=\mathrm{rc}_q-\mathrm{rc}_{\min}\ge 0$, and $\mu'$ remains non-positive. Weak
duality at this repaired point gives
\begin{equation*}
z_{\mathrm{LP}}\ge\sum_p\sigma_p+|S|\mu'=z_{\mathrm{RMP}}+|S|\cdot\mathrm{rc}_{\min}.
\end{equation*}
Substituting the pricing solve's own dual bound $\mathrm{rc}_{\mathrm{lb}}$, which never exceeds
$\mathrm{rc}_{\min}$, only weakens the inequality, and taking $\min(0,\cdot)$ leaves the bound at
$z_{\mathrm{RMP}}$ once the master has priced out. Because the covering rows relax the partition
and the partition model reproduces the visit count exactly, the chain
$\mathrm{LB}\le z_{\mathrm{LP}}\le z_{\mathrm{IP}}\le$ the optimal visit count holds.

Three conditions make the argument valid and the implementation meets all three. Pricing is solved
as an exact integer programme, so its dual bound holds even when a call is truncated. The
station-indexed repair is applied one convexity row at a time at a common dual vector, which the
run enforces after the generation loop. The bound covers the multi-item supports, to which the
singleton constant of Section~4.2.1 of the main article is added. Degeneracy of the covering master
slows how quickly the floor rises as columns arrive but never invalidates it, and we declare
convergence only when an exact solve at the true duals proves that no bundle of negative reduced
cost remains.

\section{The industrial community-bound search}\label{sup:betasearch}
Section~6 of the main article reports that the community bound $\beta$ does not transfer to
Company~A, whose stations hold between 8 and 2{,}575 movable products. This section records the
full search behind that finding.

Applying the rule $\beta=\max(5,\lfloor 0.4\,\zeta+\tfrac12\rfloor)$ to the mean station capacity
returns $\beta=266$, which no repair can make feasible. A community sized for the average station
cannot fit the small ones, so placement overloads them and the repair spends itself undoing
placement rather than balancing it.

We evaluated nineteen settings in all. Five yield a feasible layout, $\{15, 33, 47, 350, 450\}$,
and the remaining fourteen, spread from 63 to 399, do not, so feasibility is not an interval in
$\beta$. We can only offer a hypothesis for the two large feasible settings, which the data does
not let us test: once $\beta$ passes the capacity of all but the largest stations, almost every
community has to be split at placement, and the split rule that then governs the layout places
products one at a time by the same admissibility test, so the blocks it produces are sized by what
the stations can take rather than by $\beta$. On that reading the middle of the range is the worst
place to sit, because communities are too large for the small stations and too small to force the
split path throughout.

The protocol a site would run is a grid anchored on the mean capacity, $\{33,67,133,200,266,333,
399\}$, keeping the feasible layout with fewest visits. Exactly one of the seven is feasible,
$\beta=33$, and it gives the heuristic row of the main article's industrial table, whose time entry
charges the whole seven-point grid for that reason. The other twelve settings were exploratory and
are excluded from the reported cost, not because they were cheap: charging all nineteen would raise
the heuristic to roughly 1,750 seconds, still an order of magnitude below the next fastest method.
Two of those exploratory settings are better and feasible, $\beta=350$ and $\beta=450$ at
$1{,}030{,}978$ and $1{,}036{,}764$ visits, but no grid we could have specified in advance selects
them, and we report them as a limitation rather than a result. On a warehouse of interchangeable
stations none of this arises, the rule's $\beta$ being feasible on all 29 benchmark instances.

\section{Per-instance column-generation convergence record}\label{sup:cgrecord}
The released archive named in the main article's data availability statement carries, for every
column-generation run, the convergence status, the restricted-master value $z_{\mathrm{RMP}}$, the
pricing dual bound $\mathrm{rc}_{\mathrm{lb}}$ returned by the last pricing call, and the deadline
the generation loop reached. Those records support the reading of Section~4.2.4 of the main
article: no run proved that no improving bundle remained, the generation loop reached its $0.70B$
allocation on all 29 instances at all four sizes, and the last dual bound was strictly negative
everywhere, from $-696$ to $-1{,}497$ at 50 SKUs and widening to $-2.3\times10^{5}$ at 2,000 SKUs.
The quantity that pushes the reported bound down is therefore the correction
$|S|\cdot\mathrm{rc}_{\mathrm{lb}}$ rather than any weakness of the relaxation, and the binding
constraint is the number of iterations reached inside the cap.

\section{Column-generation formulation of the relocation cap}\label{sup:relocap}
The bounded-reassignment limit of Section~7.3 of the main article fits the column-generation
framework as well as the set-variable MILP. Counting, for each bundle $q$, how many of its products
would leave their legacy station,
\[
d_q \;=\; \bigl|\{\, p \in q \,:\, s'(p) \neq s(q) \,\}\bigr|,
\]
where $s'(p)$ is the legacy station of product $p$ and $s(q)$ the station bundle $q$ is installed
at, the cap becomes the single knapsack-style master row $\sum_q d_q\, y_q \le k$. Its dual enters
pricing as a per-product relocation penalty, so the subproblem proposes a bundle only when the
visits it saves outweigh the relocation budget it consumes. We have not run that variant, and the
results reported in the main article for bounded reassignment all come from the set-variable MILP.

\section{Notation}\label{sup:notation}
Table~\ref{tab:notation} collects the symbols used throughout the main article. Every symbol is
also defined at its first use there; the table is a reference aid rather than the definition.

\begin{table}[H]
\centering
\tbl{Notation used throughout the main article.}
{\begin{tabularx}{\textwidth}{@{}c@{\hspace{2em}}X@{}}
\toprule
Symbol & Definition \\
\midrule
\multicolumn{2}{@{}l}{\emph{Problem data and compact models (Sections~3 and~3.2 of the main article)}}\\
$P,\,O,\,S$ & Sets of products (SKUs), orders, and stations (zones) \\
$P_o$       & Set of products requested in order $o \in O$ \\
$\zeta_s$   & Number of storage locations at station $s$ \\
$\zeta$     & Common station capacity under homogeneous stations ($|P|/|S|$) \\
$L_p$       & Order lines of product $p$ over the planning horizon \\
$V_s$       & Processing rate of station $s$ (lines per unit time) \\
$T_s$       & Workload budget of station $s$, in the units of $\sum_p L_p/V_s$ \\
$C_s$       & Line capacity $T_s V_s$ of station $s$ over the horizon \\
$x_{ps}$    & 1 if product $p$ is assigned to station $s$ \\
$z_{os}$    & 1 if order $o$ visits station $s$ \\
$\Phi_s$    & Set variable: the subset of products assigned to station $s$ \\
\midrule
\multicolumn{2}{@{}l}{\emph{Decomposition (Section~4.2 of the main article)}}\\
$K_s$       & Feasible patterns for station $s$ (product subsets it could hold) \\
$U,\,w_u$   & Distinct multi-item order supports and their order counts \\
$Q'$        & Pool of generated bundles \\
$c_q$       & Weighted-support cost of pattern or bundle $q$ \\
$y_{sq},\,y_q$ & Selection of $q$ (station-indexed; aggregated) \\
$\sigma_p,\,\mu,\,\mu_s$ & Duals of the covering, cardinality, and convexity rows \\
$a_p,\,z_u$ & Pricing variables: membership of product $p$; touch of support $u$ \\
$\rho,\,B$  & Workload penalty rate of the swap descent; time cap of the drive \\
$\Pi,\,\Pi^{\ast},\,\Pi(s)$ & A layout; the incumbent; the products $\Pi$ places at station $s$ \\
$z_{\mathrm{RMP}},\,z_{\mathrm{IP}},\,\mathrm{rc}_{\mathrm{lb}}$ & Restricted-master value; integer-master value; pricing dual bound \\
\bottomrule
\end{tabularx}}
\label{tab:notation}
\end{table}

\bibliographystyle{apacite}
\bibliography{IJPR_CSLAP_v4}

\end{document}
